% Generated draft: requires mathematical and source review.
\documentclass[11pt]{article}
\usepackage[margin=25mm]{geometry}
\usepackage{fontspec}
\setmainfont{Latin Modern Roman}
\usepackage{amsmath,amssymb,mathtools,unicode-math}
\setmathfont{Latin Modern Math}
\usepackage{enumitem,needspace,xurl,hyperref}
\setlength{\parindent}{0pt}
\setlength{\parskip}{5pt}
\setlength{\emergencystretch}{3em}
\newcommand{\N}{\mathbb{N}}
\newcommand{\Z}{\mathbb{Z}}
\newcommand{\Q}{\mathbb{Q}}
\newcommand{\R}{\mathbb{R}}
\newcommand{\C}{\mathbb{C}}
\newcommand{\F}{\mathbb{F}}
\newcommand{\A}{\mathbb{A}}
\newcommand{\PP}{\mathbb{P}}
\newcommand{\E}{\mathbb{E}}
\newcommand{\eps}{\varepsilon}
\newcommand{\dd}{\,\mathrm{d}}
\begin{document}
\section*{UnsolvedMath ID: 10400217}
\textbf{Source problem code:} catalogue entry
\subsection{Definitions and mathematical statement}
KTG denotes knotted trivalent graphs, with the framing/vertex and graph-operation conventions of the source. The excerpt does not define their algebra, generators or relations. Construct, using that algebra, two closed-three-manifold invariants: a universal Reshetikhin--Turaev invariant and a universal Turaev--Viro--Ocneanu invariant.

A possible universality requirement is to produce values \(U_{\mathrm{RT}}(M)\) and \(U_{\mathrm{TVO}}(M)\) in specified KTG-based targets such that every admitted specialization \(\theta\) recovers its corresponding invariant,
\[
 \operatorname{ev}_\theta(U_{\mathrm{RT}}(M))=Z^{\mathrm{RT}}_\theta(M),
 \qquad
 \operatorname{ev}_\theta(U_{\mathrm{TVO}}(M))=Z^{\mathrm{TVO}}_\theta(M).
\]
This spells out a conventional interpretation, not a recovered list of allowed quantum groups, categories, roots or levels. The source does not prescribe the universal property, targets/completions, normalization or orientation/framing conventions.

The requested construction must use the intended KTG operations and be independent of the chosen manifold presentation. Merely collecting known numerical invariants is not silently identified with such a construction. The two universal theories remain distinct; no square or conjugation relation between them is assumed without hypotheses. Nor is universality for quantum specializations replaced by finite-type universality.

The summary gives only an open/triage description and no exact current construction or theorem. The 2002 program remains source information pending primary-literature review.
\subsection{Short English statement}
Build KTG-algebra universal versions of both Reshetikhin--Turaev and Turaev--Viro--Ocneanu invariants for closed three-manifolds. The KTG algebra, admissible specializations and precise universal property are not supplied.
\subsection{Sources}
\begin{itemize}
\item Local registry: \url{data/unsolvedmath/problems.json}, permanent UnsolvedMath ID \(10400217\), source code catalogue entry; supplied statement, background, and research summary, updated 2026-08-06.
\item Ohtsuki, \emph{Problems on invariants of knots and 3-manifolds} (2002), Problem 12.12, PDF page 165: \url{https://msp.org/gtm/2002/04/gtm-2002-04-024s.pdf}. Supplied PDF-font-structure extraction accessed 2026-07-29; provenance distinguishes 231 numbered items.
\item No full KTG-algebra definition, exact universal-property statement or separate modern resolution reference is supplied.
\end{itemize}
\end{document}
